% GENERATED FILE. Edit canonical lessons, then run npm run generate:books.
% canonical-source-hash: fnv1a64:ba8427e6b451e9ff
% canonical-lessons: TE-C226-patam, TE-C226-patram, TE-C226-phailu, TE-C226-pharam, TE-C226-kardu

\chapter{Map, Document, File, Form, Card}
\label{ch:te-map-document-file-form-card}
\hlchaptermodality{226}

\begin{chapteropening}
\textbf{By the end of this chapter:} \emph{I can say a map, a document, a file, a form and a card.}

Complete the last lesson of chapter 226: I can say a map, a document, a file, a form and a card.
\end{chapteropening}

\section[\texorpdfstring{\te{పటం} (paṭaṁ)}{పటం (paṭaṁ)}]{\te{పటం} (paṭaṁ) — a map, a chart}

\label{lesson:TE-C226-patam}



\begin{quote}
Before the new one: say the Telugu for a watermelon, then the Telugu for cashew nuts.
\end{quote}

\subsection*{You'll want to know: \te{పటం}}
\textbf{\te{పటం}} — \emph{paṭaṁ} — \textquotedblleft{}a map, a chart\textquotedblright{}.

\begin{grammarlens}[title={What you've built}]
One more word for this chapter.
\end{grammarlens}

\subsection*{Your turn}
\begin{itemize}
\raggedright
  \item \emph{Say it:} \emph{paṭaṁ}
  \item \emph{Say it:} \emph{paṭaṁ}, once more
  \item \emph{Say it:} say \emph{paṭaṁ}
  \item \emph{Recall:} say the Telugu for a watermelon, then the Telugu for cashew nuts, then say \emph{paṭaṁ} again
\end{itemize}

\begin{tcolorbox}[breakable,colback=teal!4,colframe=teal!35!black,title={Before you move on}]
What does \te{పటం} mean? (\textquotedblleft{}A map, a chart\textquotedblright{}.) Say it once more.
\end{tcolorbox}
\section[\texorpdfstring{\te{పత్రం} (patraṁ)}{పత్రం (patraṁ)}]{\te{పత్రం} (patraṁ) — a document, a certificate}

\label{lesson:TE-C226-patram}



\begin{quote}
Before the new one: say the Telugu for cashew nuts, then the Telugu for a map.
\end{quote}

\subsection*{You'll want to know: \te{పత్రం}}
\textbf{\te{పత్రం}} — \emph{patraṁ} — \textquotedblleft{}a document, a certificate\textquotedblright{}.

\begin{grammarlens}[title={What you've built}]
One more word for this chapter.
\end{grammarlens}

\subsection*{Your turn}
\begin{itemize}
\raggedright
  \item \emph{Say it:} \emph{patraṁ}
  \item \emph{Say it:} \emph{patraṁ}, once more
  \item \emph{Say it:} say \emph{patraṁ}
  \item \emph{Recall:} say the Telugu for cashew nuts, then the Telugu for a map, then say \emph{patraṁ} again
\end{itemize}

\begin{tcolorbox}[breakable,colback=teal!4,colframe=teal!35!black,title={Before you move on}]
What does \te{పత్రం} mean? (\textquotedblleft{}A document, a certificate\textquotedblright{}.) Say it once more.
\end{tcolorbox}
\section[\texorpdfstring{\te{ఫైలు} (phailu)}{ఫైలు (phailu)}]{\te{ఫైలు} (phailu) — a file, a folder}

\label{lesson:TE-C226-phailu}



\begin{quote}
Before the new one: say the Telugu for a map, then the Telugu for a document.
\end{quote}

\subsection*{You'll want to know: \te{ఫైలు}}
\textbf{\te{ఫైలు}} — \emph{phailu} — \textquotedblleft{}a file, a folder\textquotedblright{}.

\begin{grammarlens}[title={What you've built}]
One more word for this chapter.
\end{grammarlens}

\subsection*{Your turn}
\begin{itemize}
\raggedright
  \item \emph{Say it:} \emph{phailu}
  \item \emph{Say it:} \emph{phailu}, once more
  \item \emph{Say it:} say \emph{phailu}
  \item \emph{Recall:} say the Telugu for a map, then the Telugu for a document, then say \emph{phailu} again
\end{itemize}

\begin{tcolorbox}[breakable,colback=teal!4,colframe=teal!35!black,title={Before you move on}]
What does \te{ఫైలు} mean? (\textquotedblleft{}A file, a folder\textquotedblright{}.) Say it once more.
\end{tcolorbox}
\section[\texorpdfstring{\te{ఫారం} (phāraṁ)}{ఫారం (phāraṁ)}]{\te{ఫారం} (phāraṁ) — a form (to fill in)}

\label{lesson:TE-C226-pharam}



\begin{quote}
Before the new one: say the Telugu for a document, then the Telugu for a file.
\end{quote}

\subsection*{You'll want to know: \te{ఫారం}}
\textbf{\te{ఫారం}} — \emph{phāraṁ} — \textquotedblleft{}a form (to fill in)\textquotedblright{}.

\begin{grammarlens}[title={What you've built}]
One more word for this chapter.
\end{grammarlens}

\subsection*{Your turn}
\begin{itemize}
\raggedright
  \item \emph{Say it:} \emph{phāraṁ}
  \item \emph{Say it:} \emph{phāraṁ}, once more
  \item \emph{Say it:} say \emph{phāraṁ}
  \item \emph{Recall:} say the Telugu for a document, then the Telugu for a file, then say \emph{phāraṁ} again
\end{itemize}

\begin{tcolorbox}[breakable,colback=teal!4,colframe=teal!35!black,title={Before you move on}]
What does \te{ఫారం} mean? (\textquotedblleft{}A form (to fill in)\textquotedblright{}.) Say it once more.
\end{tcolorbox}
\section[\texorpdfstring{\te{కార్డు} (kārḍu)}{కార్డు (kārḍu)}]{\te{కార్డు} (kārḍu) — a card}

\label{lesson:TE-C226-kardu}



\begin{quote}
Before the new one: say the Telugu for a file, then the Telugu for a form.
\end{quote}

\subsection*{You'll want to know: \te{కార్డు}}
\textbf{\te{కార్డు}} — \emph{kārḍu} — \textquotedblleft{}a card\textquotedblright{}.

\begin{grammarlens}[title={What you've built}]
One more word for this chapter.
\end{grammarlens}

\subsection*{Your turn}
\begin{itemize}
\raggedright
  \item \emph{Say it:} \emph{kārḍu}
  \item \emph{Say it:} \emph{kārḍu}, once more
  \item \emph{Say it:} say \emph{kārḍu}
  \item \emph{Recall:} say the Telugu for a file, then the Telugu for a form, then say \emph{kārḍu} again
\end{itemize}

\begin{tcolorbox}[breakable,colback=teal!4,colframe=teal!35!black,title={Before you move on}]
What does \te{కార్డు} mean? (\textquotedblleft{}A card\textquotedblright{}.) Say it once more.
\end{tcolorbox}
